Recent advances in LLMs have demonstrated their potential as reading assistants. Chen and Leitch~\cite{b6} evaluated LLMs such as Claude for academic reading, reporting improvements in comprehension. However, they also highlighted limitations, including the models’ tendency to distract and concerns about privacy due to reliance on cloud infrastructure. Kryściński et al.~\cite{b7} introduced BOOKSUM, showing that LLMs can summarize long-form literary works such as novels and plays. Similarly, Spatharioti et al.~\cite{b8} found that LLM-enhanced web search improves the efficiency of information retrieval compared to traditional methods, particularly for consumer tasks. In the domain of document-level retrieval, Dense Passage Retrieval (DPR)~\cite{b9} was a significant step toward semantic search by learning dense representations of both queries and documents. However, as Reichman and Heck~\cite{b10} observed, DPR systems are typically trained on general corpora and are not easily adapted to specific, unseen documents. Improvements such as RocketQA~\cite{b11}, which leverages cross-batch negatives and distillation techniques, enhance retrieval quality but impose computational demands that limit applicability in low-resource or offline settings. Embedding models like nomic-embed-text-v2~\cite{b12}, built using a Mixture-of-Experts framework and trained on multilingual objectives, provide a balance between expressiveness and efficiency. Recent work, such as InfoCSE~\cite{b13}, has shown that masked contrastive objectives improve semantic alignment in retrieval tasks.

RAG has emerged as a robust strategy for combining semantic retrieval with generative inference. The canonical RAG architecture has been formalized and extended in works by Gao et al.~\cite{b14} and Wang et al.~\cite{b15}, including modular and task-adaptive variants. Vertical applications of RAG have been explored in domains such as medicine~\cite{b16,b17} and education~\cite{b18,b19}, typically in multi-document settings. However, user studies such as Levonian et al.~\cite{b19} indicate that retrieval strictness can conflict with user expectations, more so in educational contexts. Emerging architectures like FLARE~\cite{b20} introduce dynamic, content-aware retrieval strategies, though they often trade off control for flexibility. Privacy-preserving and local AI systems have gained attention due to increasing concerns over data security and user autonomy. Xu et al.~\cite{b21} demonstrated that on-device AI is feasible through techniques such as model quantization and runtime optimization. Lightweight infrastructure—including SQLite~\cite{b22}, Qdrant~\cite{b23}, and local inference frameworks like Ollama~\cite{b24}—supports practical deployment of vector search and inference on edge devices. However, as Mireshghallah et al.~\cite{b25} and Ippolito et al.~\cite{b26} showed, even pretrained models can leak memorized content, showing the importance of avoiding fine-tuning on user data. Similar privacy concerns were identified in on-device audio systems by Zhou et al.~\cite{b27}, who advocated for strictly offline processing to protect user interactions.

In summary, existing research has explored AI-assisted reading, semantic retrieval, RAG frameworks, and private inference independently. However, there remains a lack of integrated systems that combine progress-aware, privacy-preserving, and document-specific AI support within a single-reader context. This gap motivates the development of MeReader, which synthesizes these strands into an offline, book-constrained assistant focused on enhancing comprehension without compromising narrative continuity or user privacy.